\subsection{多项式代数上的复合算子}

设 K 为特征 0 的交换域，K{[}X{]} 为 K 上一个未定元的多项式代数（Alg., IV. 1）；在本节中，我们把 K{[}X{]} 上的算子理解为向量空间 K{[}X{]} （在 K 上）到其自身的线性映射 U；由于单项式 \(X^{n}\) （ \(n \geqslant 0\) ）构成此空间的一组基，U 由多项式 \(U(X^{n})\) 确定；确切地说，若 \(f(X) = \sum_{k=0}^{\infty} \lambda_{k} X^{k}\) ，其中 \(\lambda_{k} \in K\) ，则 \(U(f) = \sum_{k=0}^{\infty} \lambda_{k} U(X^{k})\) 。

若 \(\mathbf{G}\) 是 \(\mathbf{K}\) 上的一个含单位元的交换代数，则 \(\mathbf{G}\) -模 \(\mathrm{G}[\mathrm{X}]\) 是由把 \(\mathbf{K}\) -向量空间 \(\mathrm{K}[\mathrm{X}]\) 的标量域扩张到 \(\mathrm{G}\) 而得到的；每个算子 \(U\) （定义在 \(\mathrm{K}[\mathrm{X}]\) 上）都以唯一的方式扩张为 \(\mathrm{G}\) -模 \(\mathrm{G}[\mathrm{X}]\) 到其自身的一个线性映射，我们仍将它记作 \(U\) （Alg., II, p.~278）；对每个元素 \(g(\mathrm{X}) = \sum_{k=0}^{\infty} \gamma_k X^k\) ，其中 \(\gamma_k \in \mathrm{G}\) ，有 \(U(g) = \sum_{k=0}^{\infty} \gamma_k U(X^k)\) 。

特别地，考虑 \(\mathrm{G} = \mathrm{K}[\mathrm{Y}]\) 的情形；这时 \(\mathrm{G}[\mathrm{X}]\) 是环 \(\mathrm{K}[\mathrm{X},\mathrm{Y}]\) ，即 \(\mathrm{K}\) 上两个未定元的多项式环；为避免混淆，我们把 \(U\) 到 \(\mathrm{G}[\mathrm{X}]\) 上的扩张记作 \(U_{\mathrm{X}}\) 。于是，对任意多项式 \(g(\mathrm{X},\mathrm{Y}) = \sum_{k=0}^{\infty}\gamma_k(\mathrm{Y})\mathrm{X}^k\) （其中 \(\gamma_k(Y)\in \mathrm{K}[\mathrm{Y}]\) ）有 \(U_{\mathrm{X}}(g) = \sum_{k=0}^{\infty}\gamma_k(\mathrm{Y})U(\mathrm{X}^k)\) 。由于 \(U_{\mathrm{X}}\) 是线性的，可以看出，若把它写成 \(g(\mathrm{X},\mathrm{Y}) = \sum_{h=0}^{\infty}\beta_h(\mathrm{X})\mathrm{Y}^h\) ，则也有 \(U_{\mathrm{X}}(g) = \sum_{h=0}^{\infty}U(\beta_h)\mathrm{Y}^h\) 。

在把 X 对应到 Y 的 K{[}X{]} 到 K{[}Y{]} 上的典范同构下，算子 U 变换为 K{[}Y{]} 上的一个算子；为避免混淆，我们把它记作 \(U_{Y}\) ，于是 \(U_{\mathrm{Y}}(f(\mathrm{Y}))\) 就是在多项式 \(U(f(\mathrm{X})) = U_{\mathrm{X}}(f(\mathrm{X}))\) 中把 X 换成 Y 所得到的多项式。这个算子 \(U_{Y}\) 又可以扩张为一个算子（仍记作 \(U_{Y}\) ），它定义在 K{[}X,Y{]} 上：若 \(g(\mathrm{X}, \mathrm{Y}) = \sum_{h=0}^{\infty} \beta_{h}(\mathrm{X}) \mathrm{Y}^{h}\) ，则 \(U_{\mathrm{Y}}(g(\mathrm{X}, \mathrm{Y})) = \sum_{h=0}^{\infty} \beta_{h}(\mathrm{X}) U_{\mathrm{Y}}(\mathrm{Y}^{h})\) 。

作为这些扩张的一个例子，我们举出 K{[}X{]} 上的求导算子 D （Alg., IV. 6），它给出偏求导算子 \(D_{X}\) 与 \(D_{Y}\) （定义在 K{[}X,Y{]} 上）。

对每个多项式 \(f \in K[X]\) ，我们用 \(T_{Y}(f)\) 表示多项式 \(f(X + Y)\) （它属于 K{[}X,Y{]}）；映射 \(T_{Y}\) 是从 K{[}X{]} 到 K{[}X,Y{]} 中的一个 K-线性映射，称为平移算子。

定义 1. 称 K{[}X{]} 上的算子 U 为复合算子，如果它与平移算子可交换，也就是说，如果 \(U_{X}T_{Y} = T_{Y}U\) 。

换句话说，若 \(f\) 是 \(\mathbf{K}[\mathbf{X}]\) 中的任意一个多项式，且 \(g = U(f)\) ，则必须有 \(g(\mathbf{X} + \mathbf{Y}) = U_{\mathbf{X}}(f(\mathbf{X} + \mathbf{Y}))\) 。

由这个定义立即得出，对每个多项式 \(f(\mathrm{X}) \in \mathrm{K}[\mathrm{X}]\) ，按上述记号有

\[
U _ {\mathrm{X}} \big (f (\mathrm{X} + \mathrm{Y}) \big) = U _ {\mathrm{Y}} \big (f (\mathrm{X} + \mathrm{Y}) \big).\tag{1}
\]

例. 1) 对每个 \(\lambda \in \mathbf{K}\) ，把每个多项式 \(f(\mathbf{X})\) 对应到多项式 \(f(\mathbf{X} + \lambda)\) 的算子是一个复合算子。

\begin{enumerate}
\def\labelenumi{\arabic{enumi})}
\setcounter{enumi}{1}
\tightlist
\item
  K{[}X{]} 上的求导 D 是一个复合算子（参见命题 1）。
\end{enumerate}

注. 由于 K 是无限域，K{[}X{]} 上的算子 U 是复合算子，当且仅当对每个多项式 \(f \in K[X]\) 和每个元素 \(\alpha \in K\) ，令 \(g = U(f)\) ，有 \(g(X + \alpha) = U(f(X + \alpha))\) 。（Alg., IV. 16, cor.）

显然，复合算子的以 K 为系数的每个线性组合仍是复合算子；两个复合算子的复合（乘积）也是如此。换句话说，复合算子构成一个子代数 \(\Gamma\) ，含于向量空间 K{[}X{]} 的自同态代数中。

命题 1. K{[}X{]} 上的算子 U 是复合算子的充分必要条件为：它与 K{[}X{]} 上的求导 D 可交换。

事实上，泰勒公式表明，对每个多项式 \(f \in K[X]\) 有

\[
U _ {\mathrm{X}} \big (f (\mathrm{X} + \mathrm{Y}) \big) = U _ {\mathrm{X}} \left(\sum_ {k = 0} ^ {\infty} \frac {1}{k !} \mathrm{Y} ^ {k} \mathrm{D} ^ {k} f (\mathrm{X})\right) = \sum_ {k = 0} ^ {\infty} \frac {1}{k !} \mathrm{Y} ^ {k} U \big (\mathrm{D} ^ {k} f (\mathrm{X}) \big);
\]

若令 \(g = U(f)\) ，则有

\[
g (\mathrm{X} + \mathrm{Y}) = \sum_ {k = 0} ^ {\infty} \frac {1}{k !} \mathrm{Y} ^ {k} \mathrm{D} ^ {k} g (\mathrm{X}) = \sum_ {k = 0} ^ {\infty} \frac {1}{k !} \mathrm{Y} ^ {k} \mathrm{D} ^ {k} (U (f (\mathrm{X})))
\]

于是，为使 U 是复合算子，必须有 \(UD^{k} = D^{k}U\) （对每个整数 \(k \geqslant 1\) ），特别地有 \(UD = DU\) 。反之，若这一关系成立，则由对 k 的归纳法，它蕴含 \(UD^{k} = D^{k}U\) （对每个整数 \(k \geqslant 1\) ）；于是泰勒公式表明 \(g(X + Y) = U_{X}(f(X + Y))\) 。

对每个多项式 \(f \in K[X, Y]\) ，我们用 \(U_{0}(f)\) 表示多项式 \(U_{X}(f)\) 中不含 X 的项；特别地，若 \(f \in K[X]\) ，则 \(U_{0}(f)\) 是 \(U(f)\) 中的常数项，而 \(U_{0}\) 是 K{[}X{]} 上的一个线性型。对每个多项式 \(f \in K[X]\) ，令 \(g = U(f)\) ；由 VI, p.~270 的定义 1

\[
g (\mathrm{X} + \mathrm{Y}) = U _ {\mathrm{X}} (f (\mathrm{X} + \mathrm{Y})) = U _ {\mathrm{X}} \left(\sum_ {k = 0} ^ {\infty} \frac {1}{k !} \mathrm{X} ^ {k} \mathrm{D} ^ {k} f (\mathrm{Y})\right) = \sum_ {k = 0} ^ {\infty} \frac {1}{k !} U (\mathrm{X} ^ {k}) \mathrm{D} ^ {k} f (\mathrm{Y})
\]

而若在这个公式中把 X 换成 0 ，就得到

\[
g (\mathrm{Y}) = \sum_ {k = 0} ^ {\infty} \frac {1}{k !} U _ {0} (\mathrm{X} ^ {k}) \mathrm{D} ^ {k} f (\mathrm{Y}).
\]

于是可以看出

\[
U (f (\mathrm{X})) = \sum_ {k = 0} ^ {\infty} \frac {1}{k !} \mu_ {k} \mathrm{D} ^ {k} f (\mathrm{X})\tag{2}
\]

其中 \(\mu_{k}\) 是多项式 \(U(\mathbf{X}^k)\) 的常数项。

这个公式表明，诸 \(\mu_{k}\) 完全确定了复合算子 U；反过来，若 \((\mu_{n})\) 是 K 中元素的一个任意序列，则公式 (2) 定义了一个算子 U，它显然与 D 可交换，因而（VI, p.~270, prop. 1）是一个复合算子。今后我们把 (2) 写成如下形式

\[
U = \sum_ {k = 0} ^ {\infty} \frac {1}{k !} \mu_ {k} D ^ {k}.\tag{3}
\]

这个公式可以用拓扑语言解释如下：若在 K{[}X{]} 上考虑离散拓扑，并在自同态代数 \(\operatorname{End}(\operatorname{K}[X])\) （由 K{[}X{]} 的自同态组成）上考虑单收敛拓扑（Gen.~Top., X, p.~277），则以 \(\frac{1}{k!}\mu_{k}D^{k}\) 为通项的级数在 \(\operatorname{End}(\operatorname{K}[X])\) 中交换收敛，且其和为 U （Gen.~Top., III, p.~269）。

公式 (3) 表明，对形式级数 \(u(S) = \sum_{k=0}^{\infty} \alpha_k S^k\) （\(K\) 上一个未定元的，Alg., IV. 41），可以对应一个复合算子 \(U = \sum_{k=0}^{\infty} \alpha_k D^k\) ，今后我们把它记作 \(u(D)\) 。这一注记可以按如下方式加以阐明：

定理 1. 把每个形式级数 \(u(\mathbf{S}) = \sum_{k=0}^{\infty} \alpha_k \mathbf{S}^k\) （\(\mathbf{K}\) 上一个未定元的）对应到复合算子 \(u(\mathbf{D}) = \sum_{k=0}^{\infty} \alpha_k \mathbf{D}^k\) （\(\mathrm{K}[\mathrm{X}]\) 上的）的映射，是形式级数代数 \(\mathrm{K}[[\mathrm{S}]]\) 到复合算子代数 \(\Gamma\) 上的一个同构。

立即可以验证这个映射是一个同态。剩下要证的是它是单射，换句话说，关系 \(\sum_{k=0}^{\infty}\alpha_{k}D^{k}=0\) 蕴含对每个 k 都有 \(\alpha_{k}=0\) ；而 \(h!\alpha_{h}\) 是把 \(\sum_{k=0}^{\infty}\alpha_{k}D^{k}\) 作用于 \(X^{h}\) 所得多项式的常数项，由此即得定理。

推论. 复合算子代数 \(\Gamma\) （K{[}X{]} 中的）是交换的。

例. 若 \(U\) 是把每个多项式 \(f(\mathrm{X})\) 对应到 \(f(\mathrm{X} + \lambda)\) （其中 \(\lambda \in \mathbf{K}\) ）的算子，则有 \(U_0(\mathrm{X}^k) = \lambda^k\) ，从而 \(U = \sum_{k=0}^{\infty} \frac{1}{k!} (\lambda \mathrm{D})^k\) 。与 \(e^x\) 的级数展开（III, p.~105）相类比，我们用 \(e^{\mathrm{S}}\) 或 \(\exp(\mathrm{S})\) 表示形式级数 \(\sum_{n=0}^{\infty} \frac{1}{n!} \mathrm{S}^n\) （环 \(\mathrm{K}[[\mathrm{S}]]\) 中的）；于是可以写 \(U = e^{\lambda \mathrm{D}}\) 。在这论证中把域 \(\mathrm{K}\) 换成有理分式域 \(\mathrm{K}(\mathrm{Y})\) ，同样可以看出平移算子 \(T_{\mathrm{Y}}\) 可写成 \(e^{\mathrm{YD}}\) 。

此外，在 K 上两个未定元的形式级数环 K{[}{[}S, T{]}{]} 中有

\[
\begin{array}{l} (\exp S) (\exp T) = \sum_ {p, q} \frac {S ^ {p}}{p !} \frac {T ^ {q}}{q !} \\ = \sum_ {n = 0} ^ {\infty} \frac {1}{n !} \left(\binom {n} {0} S ^ {n} + \binom {n} {1} S ^ {n - 1} T + \dots + \binom {n} {n} T ^ {n}\right) \\ = \exp (S + T) \end{array}\tag{4}
\]

特别地，有

\[
(\exp S) (\exp (- S)) = 1\tag{5}
\]

这就说明我们引入的记号是合理的。

按语. 形式级数代数 K{[}{[}S{]}{]} 与复合算子代数 \(\Gamma\) （K{[}X{]} 上的）之间的同构，有时使我们能用对相应的复合算子证明命题的办法，最简单地证明关于形式级数的命题（参见 VI, p.~274, prop. 6）。

